% Asset ID: AS1AlgebraicExpressionsTIKZ-005
% Source: Dr Frost fractional and negative indices examples
% Related lesson section: Negative and Fractional Indices, Worked Examples
% Purpose: Show how a negative fractional power is unpacked using reciprocal, root and power.

\documentclass[tikz,border=8pt]{standalone}
\usetikzlibrary{arrows.meta,positioning}

\begin{document}
\begin{tikzpicture}[
    font=\sffamily,
    arrow/.style={-{Latex[length=3mm]}, thick},
    step/.style={draw, rounded corners=4pt, thick, fill=gray!8, align=center, text width=4.4cm, minimum height=1.2cm},
    final/.style={draw, rounded corners=4pt, thick, fill=black, text=white, align=center, text width=4.4cm, minimum height=1.2cm}
]

\node[font=\sffamily\bfseries\Large] at (0,3.2)
{Unpacking a negative fractional index};

\node[step] (start) at (0,2.0)
{$\left(\dfrac{27}{8}\right)^{-2/3}$};

\node[step] (recip) at (-4,0.35)
{Negative sign\\take reciprocal\\$\left(\dfrac{8}{27}\right)^{2/3}$};

\node[step] (root) at (0,-1.35)
{Denominator $3$\\take cube root\\$\left(\dfrac{2}{3}\right)^2$};

\node[step] (power) at (4,0.35)
{Numerator $2$\\square the result};

\node[final] (answer) at (0,-3.05)
{$\left(\dfrac23\right)^2=\dfrac49$};

\draw[arrow] (start.south west) -- (recip.north);
\draw[arrow] (recip.south east) -- (root.north west);
\draw[arrow] (root.north east) -- (power.south west);
\draw[arrow] (root.south) -- (answer.north);

\node[draw, rounded corners=4pt, thick, fill=gray!12, align=center, text width=9.5cm] at (0,-4.45)
{\textbf{Memory line:} negative means reciprocal, denominator means root, numerator means power.};

\end{tikzpicture}
\end{document}
